% !TeX root = ../tcc.tex
\chapter{Planejamento do estudo utilizando ciência de dados}
\label{cap:metodologia-experimental}
\begingroup
% Mantém o espaçamento natural quando uma tabela passa para a página seguinte.
\raggedbottom

Este capítulo planeja um estudo de ciência de dados que avaliará dois
classificadores binários de fotografias clínicas de pele, implementados com
ResNet-50 e MobileNetV2. O estudo empregará dados secundários de uma
coleção pública, obtidos por pesquisa em arquivo, com partições separadas
por grupos de fotografias, e os dois classificadores serão avaliados nas
mesmas fotografias do teste reservado.

O planejamento responde às lacunas identificadas no
Capítulo~\ref{cap:revisao-literatura}: apenas 2 dos 16 estudos que
desenvolvem modelos declaram separar os dados por paciente, só um testa o
modelo em outra base, nenhum dos 17 estudos com dados reporta desempenho
por tom de pele, e os estudos que isolam a hanseníase relatam
sensibilidade bem inferior aos indicadores globais. A disponibilidade de
coleções públicas com imagens de hanseníase e de outras dermatoses e de
redes pré-treinadas torna viável examinar essas questões com um protocolo
reprodutível. Como os dados já existem e a análise é preditiva, o estudo
segue o padrão \textit{Data Science} da ACM SIGSOFT
\cite{acmSigsoftDataScience}. A definição do objetivo e a análise das
ameaças à validade se apoiam nos capítulos 8 e 9 de \citeonline{wohlin2024}.

O padrão organiza seus critérios em atributos essenciais, que todo estudo
do tipo deve atender, e atributos desejáveis, que fortalecem o estudo sem
ser obrigatórios. A Tabela~\ref{tab:atributos-data-science} relaciona cada
atributo essencial às partes do trabalho: a coluna central indica onde os
capítulos anteriores e os apêndices fundamentam ou motivam o atributo, e a
coluna da direita, onde este capítulo o atende. Entre os atributos
desejáveis, o planejamento adota o pacote de replicação, o uso de mais de
um tipo de modelo, a repetição do treinamento com partições e sementes
diferentes, a seleção criteriosa de hiperparâmetros e a inspeção manual de
parte dos dados.

\begingroup
\footnotesize
\setlength{\tabcolsep}{4pt}
\begin{longtable}{|>{\raggedright\arraybackslash}p{3.3cm}|>{\raggedright\arraybackslash}p{6.0cm}|>{\raggedright\arraybackslash}p{5.6cm}|}
	\caption{Atributos essenciais do padrão \textit{Data Science} e sua relação com as partes do trabalho}\label{tab:atributos-data-science} \\
	\hline
	\textbf{Atributo essencial} & \textbf{Fundamentação nos Capítulos~\ref{cap:introducao} a~\ref{cap:revisao-literatura} e nos apêndices} & \textbf{Atendimento no Capítulo~\ref{cap:metodologia-experimental}} \\ \hline
	\endfirsthead
	\multicolumn{3}{l}{\footnotesize\textit{Tabela~\ref{tab:atributos-data-science} (continuação)}} \\
	\hline
	\textbf{Atributo essencial} & \textbf{Fundamentação nos Capítulos~\ref{cap:introducao} a~\ref{cap:revisao-literatura} e nos apêndices} & \textbf{Atendimento no Capítulo~\ref{cap:metodologia-experimental}} \\ \hline
	\endhead
	\multicolumn{3}{r}{\footnotesize\textit{continua na próxima página}} \\
	\endfoot
	\endlastfoot
	Por que é oportuno investigar o problema com o método proposto
	& Seções~\ref{contexto} e~\ref{problema}: relevância da hanseníase e problema de pesquisa. Subseções~\ref{subsec:rq3-metodos}, \ref{subsec:rq5-rq6-tom-pele} e~\ref{subsec:rq7-generalizacao}: lacunas de separação por paciente, de resultados por tom de pele e de teste em outras bases
	& Introdução do capítulo: lacunas que motivam o estudo e disponibilidade de coleções públicas e de redes pré-treinadas \\ \hline
	Como e por que os dados foram selecionados
	& Subseção~\ref{subsec:rq2-datasets} e Apêndice~\ref{apend:inventario-datasets}: conjuntos usados na literatura, origem e condições de acesso
	& Subseção~\ref{subsec:justificativa-colecao}: escolha do 32 CSD v2 e descarte do CO2Wounds-V2 e do AI4Leprosy \\ \hline
	Configuração experimental, por exemplo com diagrama de fluxo de dados
	& Seção~\ref{gqm-adaptado-e-questuxe3o-de-pesquisa}: GQM. Subseção~\ref{subsec:ciencia-de-dados}: papéis da validação e do teste reservado
	& Seção~\ref{sec:objetivo-avaliacao-empirica}; Seção~\ref{sec:configuracao-avaliacao-empirica}, com o diagrama de fluxo de dados (Figura~\ref{fig:dfd-estudo}) e o diagrama de atividades (Figura~\ref{fig:atividades-estudo}) \\ \hline
	Engenharia de atributos e transformações aplicadas
	& Seção~\ref{sec:ia-visao-computacional}: representações aprendidas por redes convolucionais. Subseção~\ref{subsec:rq1-tecnicas}: preparo das imagens, aumento de dados e atributos usados nos estudos
	& Subseção~\ref{subsec:preprocessamento-arquiteturas}: redimensionamento, normalização e atributos das camadas convolucionais; aumento de dados restrito ao treino \\ \hline
	Pré-processamento, filtragem e categorização dos dados
	& Subseção~\ref{subsec:ciencia-de-dados}: unidade de agrupamento e vazamento entre treino e teste. Subseções~\ref{subsec:rq2-datasets} e~\ref{subsec:rq3-metodos}: unidades de amostra e divisão por imagem na literatura
	& Subseção~\ref{subsec:curadoria-procedencia}: seleção, harmonização dos rótulos e agrupamento de duplicatas. Subseção~\ref{subsec:particoes}: divisão por grupos \\ \hline
	Linhas de base do estado da arte, ou justificativa de sua ausência ou da impraticabilidade da comparação direta
	& Subseção~\ref{subsec:ciencia-de-dados}: exigência do padrão. Subseções~\ref{subsec:rq4-metricas} e~\ref{subsec:rq6-modelos} e Apêndice~\ref{apend:inventario-modelos}: modelos e resultados relatados nos estudos
	& Subseção~\ref{subsec:linhas-base}: discussão das três alternativas do padrão, classificador trivial e regressão com metadados \\ \hline
	Definição da abordagem de modelagem, em geral com pseudocódigo
	& Seção~\ref{sec:ia-visao-computacional}: aprendizado profundo e tarefas de análise de imagens. Subseção~\ref{subsec:rq6-modelos} e Apêndice~\ref{apend:inventario-modelos}: uso da ResNet-50 e da MobileNetV2 nos estudos
	& Subseções~\ref{subsec:preprocessamento-arquiteturas} a~\ref{subsec:hiperparametros} e Algoritmo~\ref{alg:procedimento-experimental} \\ \hline
	Infraestrutura de hardware e software
	& Subseção~\ref{subsec:ciencia-de-dados}: disponibilização de materiais para replicação
	& Seção~\ref{sec:instrumentacao}: objetos, diretrizes, instrumentos de medição, equipamento e bibliotecas \\ \hline
	Justificativa das estatísticas e das heurísticas
	& Subseção~\ref{subsec:ciencia-de-dados}: incerteza das estimativas. Seção~\ref{sec:avaliacao-generalizacao}: interpretação conforme o desenho do estudo
	& Subseções~\ref{subsec:curadoria-procedencia} (limiares de similaridade), \ref{subsec:particoes} (tamanho do teste), \ref{subsec:procedimento-validacao} (validação cruzada e índice de Youden) e~\ref{subsec:hiperparametros} (busca em grade); Seção~\ref{sec:medidas-avaliacao-empirica} (\textit{bootstrap} de grupos) \\ \hline
	Descrição e justificativa das métricas
	& Subseção~\ref{subsec:metricas-avaliacao}: definição das métricas \metrica{M1} a \metrica{M12}. Subseção~\ref{subsec:rq4-metricas}: métricas relatadas nos estudos
	& Tabela~\ref{tab:gqm-questoes-metricas} e Seção~\ref{sec:medidas-avaliacao-empirica}: acurácia balanceada como medida principal e medidas complementares \\ \hline
	Resultados além de valores únicos (variação e incerteza)
	& Subseção~\ref{subsec:ciencia-de-dados}: relato da variabilidade. Subseção~\ref{subsec:rq4-metricas}: forma como os estudos relatam os resultados
	& Seção~\ref{sec:medidas-avaliacao-empirica}: intervalos de 95\%, diferença pareada e variação entre dobras e sementes \\ \hline
	Pressupostos técnicos e ameaças à validade específicos de ciência de dados
	& Subseção~\ref{subsec:ciencia-de-dados} e Seção~\ref{sec:avaliacao-generalizacao}: vazamento, avaliação interna e externa. Subseções~\ref{subsec:rq5-rq6-tom-pele} e~\ref{subsec:rq7-generalizacao}: limites das evidências por tom de pele e em outras bases
	& Seção~\ref{sec:validade-avaliacao-empirica}: pressupostos técnicos e validade de construto, interna, de conclusão e externa \\ \hline
\end{longtable}
\endgroup
\fonte{Autores.}

\section{Objetivo e questões do estudo}
\label{sec:objetivo-avaliacao-empirica}

\begin{quote}
	Qual é o desempenho dos classificadores baseados em ResNet-50 e MobileNetV2
	para distinguir fotografias anotadas como hanseníase de fotografias anotadas
	como outras condições de pele, quando avaliados no mesmo conjunto de teste
	reservado?
\end{quote}

O objetivo segue os cinco campos do GQM apresentado na
Seção~\ref{gqm-adaptado-e-questuxe3o-de-pesquisa}, em consonância com a
definição de escopo da Seção~8.1 de \citeonline{wohlin2024}
(Tabela~\ref{tab:gqm-avaliacao-empirica}). A unidade de predição e
avaliação será a fotografia, e os rótulos da coleção serão a referência
para o desempenho.

\begingroup
\begin{longtable}{|p{0.19\linewidth}|p{0.75\linewidth}|}
	\caption{Objetivo do estudo segundo o GQM}\label{tab:gqm-avaliacao-empirica} \\
	\hline
	\textbf{Campo} & \textbf{Descrição} \\ \hline
	\endfirsthead
	\multicolumn{2}{l}{\footnotesize\textit{Tabela~\ref{tab:gqm-avaliacao-empirica} (continuação)}} \\
	\hline
	\textbf{Campo} & \textbf{Descrição} \\ \hline
	\endhead
	\multicolumn{2}{r}{\footnotesize\textit{continua na próxima página}} \\
	\endfoot
	\endlastfoot
	Objeto      & Classificadores binários baseados em ResNet-50 e MobileNetV2, aplicados a fotografias clínicas de pele.                 \\ \hline
	Propósito   & Estimar e comparar o desempenho dos dois classificadores no mesmo teste reservado.                                      \\ \hline
	Foco        & Distinção entre fotografias anotadas como hanseníase e fotografias anotadas como outras condições de pele.              \\ \hline
	Perspectiva & Pesquisadores de engenharia de software e da área da saúde interessados na avaliação de modelos de visão computacional. \\ \hline
	Contexto    & Estudo computacional retrospectivo com dados públicos. O apoio à triagem no SUS motiva o estudo.                        \\ \hline
\end{longtable}
\endgroup
\fonte{Autores.}

Seguindo o GQM, o objetivo se desdobra em questões, e cada questão é
respondida por métricas definidas na Subseção~\ref{subsec:metricas-avaliacao}.
Ao longo deste capítulo, cada métrica é citada pelo nome, e seu código é
uma hiperligação para essa definição.
A Tabela~\ref{tab:gqm-questoes-metricas} conecta essas métricas ao seu uso
na literatura revisada (Tabela~\ref{tab:metricas-codigos}). A variação de desempenho entre bases
(\metrica{M7}), a sobreposição espacial (\metrica{M8}) e a precisão média
na detecção (\metrica{M9}) não se aplicam, porque o estudo usa uma única
coleção e a tarefa é de classificação.

\begingroup
\begin{longtable}{|>{\footnotesize\raggedright\arraybackslash}p{4.6cm}|>{\footnotesize\raggedright\arraybackslash}p{5.8cm}|>{\footnotesize\raggedright\arraybackslash}p{4.0cm}|}
	\caption{Questões e métricas do GQM do estudo e seu uso na literatura revisada}\label{tab:gqm-questoes-metricas} \\
	\hline
	\textbf{Questão} & \textbf{Métricas} & \textbf{Uso na literatura (n/17)} \\ \hline
	\endfirsthead
	\multicolumn{3}{l}{\footnotesize\textit{Tabela~\ref{tab:gqm-questoes-metricas} (continuação)}} \\
	\hline
	\textbf{Questão} & \textbf{Métricas} & \textbf{Uso na literatura (n/17)} \\ \hline
	\endhead
	\multicolumn{3}{r}{\footnotesize\textit{continua na próxima página}} \\
	\endfoot
	\endlastfoot
	Q1. Qual é a capacidade de cada classificador de distinguir fotografias de hanseníase das demais?
	& Acurácia balanceada (\metrica{M11}, principal), sensibilidade (\metrica{M1}), especificidade (\metrica{M2}) e área sob a curva ROC (\metrica{M10}), com intervalos de 95\%
	& Acurácia balanceada: 0; sensibilidade: 15; especificidade: 6; área sob a curva ROC: 4 \\ \hline
	Q2. O desempenho difere entre ResNet-50 e MobileNetV2?
	& Diferença pareada da acurácia balanceada (\metrica{M11}), da sensibilidade (\metrica{M1}) e da especificidade (\metrica{M2}), com intervalo de 95\%
	& Não há comparação pareada no mesmo teste \\ \hline
	Q3. Os resultados podem ser situados em relação aos relatados na literatura?
	& Acurácia (\metrica{M5}), precisão (\metrica{M3}) e $F_1$ (\metrica{M4}), relatadas para discussão qualitativa
	& Acurácia: 15; precisão: 10; $F_1$: 10 \\ \hline
	Q4. Os escores dos classificadores estão calibrados?
	& Escore de Brier e erro esperado de calibração (\metrica{M12}), com diagrama de confiabilidade
	& Brier ou erro esperado de calibração: 0 \\ \hline
	Q5. O desempenho varia entre faixas de tom de pele?
	& Diferença de desempenho entre grupos (\metrica{M6}) na sensibilidade, na especificidade e na acurácia balanceada, entre faixas de ITA, de forma exploratória
	& Diferença entre grupos de tom de pele: 0 \\ \hline
	Q6. Atributos da fonte bastam para prever a classe?
	& Acurácia balanceada (\metrica{M11}) da linha de base com metadados, comparada à dos classificadores
	& Não avaliado nos estudos \\ \hline
\end{longtable}
\endgroup
\fonte{Autores.}

\section{Seleção, curadoria e particionamento dos dados}
\label{sec:particionamento-avaliacao-empirica}

\subsection{Justificativa da coleção selecionada}
\label{subsec:justificativa-colecao}

O estudo utilizará o \textit{32 Curated Categories of Skin Disease
	Images} (32 CSD), versão 2, identificado em busca complementar por dados
públicos. O arquivo distribuído nessa versão contém 5.034 imagens em 32
categorias, das quais 672 estão distribuídas entre três categorias de
hanseníase. A coleção reúne imagens de \textit{31 Classes of Skin Disease}
e \textit{Our Normal Skin} e, por meio do primeiro conjunto, incorpora
imagens do Atlas Dermatológico e do ISIC 2019
\cite{cahyantoDados2026,atlasHanseniaseDimorfa,isicChallenge2019}. Ela foi
escolhida por ser de acesso aberto, por conter fotografias de hanseníase e
de várias outras dermatoses, o que permite formar as duas classes na mesma
coleção, e por documentar a origem de parte das imagens.

As demais coleções identificadas não atendem a esse desenho. O
CO2Wounds-V2 contém fotografias de feridas crônicas em pessoas com
hanseníase, com máscaras da região da ferida, e carece de exemplos de
outras condições de pele \cite{co2woundsDados2024}. O AI4Leprosy tem
acesso condicionado a cadastro e validação \cite{barbieri2022}, que não foi
concedido pelos responsáveis.

A origem comum das imagens exige cautela. Uma verificação preliminar
encontrou a mesma fotografia de hanseníase dimorfa do Atlas Dermatológico
no 32 CSD e no \textit{Fredrick Skin Disease Dataset}
\cite{atlasHanseniaseDimorfa,fredrickSkinDisease2025}, usado por
\citeonline{eita2025}, e \citeonline{lima2024skin} também empregaram imagens
do Atlas. Por isso, os resultados deste estudo não serão comparados
numericamente com os da literatura como se viessem de dados independentes,
e a procedência será auditada na curadoria.

\subsection{Curadoria e auditoria da procedência}
\label{subsec:curadoria-procedencia}

O inventário registrará a versão, o nome do arquivo, o identificador
SHA-256, a modalidade, o rótulo original, o rótulo binário, a procedência
conhecida, o grupo de fotografias relacionadas e o motivo de exclusão,
quando houver. A curadoria ocorrerá nas seguintes etapas.

\begin{enumerate}
	\item \textbf{Selecionar fotografias e harmonizar rótulos.} Serão
	      consideradas fotografias clínicas RGB. As três categorias de
	      hanseníase formarão a classe positiva e as demais condições de pele
	      elegíveis, a classe negativa. As imagens
	      de dermatoscopia e os arquivos corrompidos ficarão fora da análise
	      principal. Dois revisores inspecionarão uma amostra por categoria e
	      fonte conhecida, registrarão divergências e documentarão casos
	      excluídos por modalidade ou rótulo inconclusivo.
	\item \textbf{Agrupar fotografias relacionadas.} Hashes criptográficos
	      localizarão cópias exatas, e a similaridade perceptual levantará
	      candidatos a duplicatas transformadas. Os limiares de similaridade
	      serão calibrados por inspeção visual de uma amostra de pares e
	      registrados no inventário. Candidatos a recortes ou outras versões
	      serão verificados por correspondência local, consistência geométrica
	      e inspeção visual. Cópias exatas com rótulos compatíveis serão
	      reduzidas a um representante, e conflitos binários sem resolução
	      serão excluídos. Fotos derivadas e registros da mesma lesão ou
	      pessoa, quando identificáveis, serão mantidos no mesmo grupo
	      \cite{kapoor2023}.
	\item \textbf{Investigar a associação entre fonte e classe.} Será
	      examinado o cruzamento entre procedência documentada e rótulo, bem
	      como diferenças de formato, bordas e marcas visuais. Se houver
	      concentração de rótulos por fonte, será verificada a disponibilidade
	      de fotografias de ambas as classes nas mesmas fontes. O controle
	      com atributos do arquivo, descrito na
	      Subseção~\ref{subsec:linhas-base}, indicará se esses atributos
	      bastam para prever a classe.
\end{enumerate}

\subsection{Partições e unidade de agrupamento}
\label{subsec:particoes}

Após a curadoria, uma divisão estratificada por classe e agrupada por
fotografias relacionadas reservará grupos distintos para desenvolvimento e
teste. Prevê-se reservar cerca de 20\% a 25\% dos dados elegíveis para
teste, e a meta de planejamento é reunir cerca de 100 grupos positivos no
teste, se a coleção curada permitir. Com 100 grupos positivos e
sensibilidade próxima de 0,80, o intervalo de confiança de 95\% teria
semiamplitude de cerca de 8 pontos percentuais, precisão considerada
suficiente para distinguir desempenhos como os relatados na revisão. A
proporção final será escolhida antes do treino, considerando a quantidade
de grupos positivos e a distribuição das condições da classe negativa.

O desenvolvimento será então dividido, também por grupos e com
estratificação por classe, em duas partes: cerca de 80\% dos grupos para a
validação cruzada e cerca de 20\% para a definição dos limiares. Serão
relatados os totais de fotografias e de grupos por partição e por classe.
As listas de grupos de cada partição serão geradas com semente fixa,
salvas antes de qualquer treinamento e não serão alteradas depois. O teste
ficará reservado até a fixação dos classificadores e de seus limiares.

\section{Configuração experimental e modelagem}
\label{sec:configuracao-avaliacao-empirica}

A configuração experimental é apresentada em duas vistas. O diagrama de
fluxo de dados da Figura~\ref{fig:dfd-estudo} mostra que dados entram em
cada processo, que dados ele produz e onde ficam armazenados, como sugere o
padrão \textit{Data Science} para descrever a configuração
\cite{acmSigsoftDataScience}. O diagrama de atividades da
Figura~\ref{fig:atividades-estudo} mostra a ordem das etapas, os pontos de
decisão e as atividades executadas em paralelo para as duas redes. O
Algoritmo~\ref{alg:procedimento-experimental} detalha o procedimento em
pseudocódigo.

\begin{figure}[H]
	\centering
	\caption{Diagrama de fluxo de dados do estudo}
	\label{fig:dfd-estudo}
	\begin{tikzpicture}[
		font=\scriptsize,
		execute at begin node={\hyphenpenalty=10000},
		entidade/.style={draw, thick, fill=black!6, align=center, text width=2.4cm, minimum height=9mm},
		processo/.style={draw, rounded corners=6pt, align=center, text width=3.0cm, minimum height=10mm},
		deposito/.style={align=left, text width=3.0cm, minimum height=7mm, inner xsep=4pt,
				append after command={\pgfextra{\draw (\tikzlastnode.north west) -- (\tikzlastnode.north east)
							(\tikzlastnode.south west) -- (\tikzlastnode.south east)
							(\tikzlastnode.north west) -- (\tikzlastnode.south west);}}},
		fluxo/.style={-{Latex[length=1.8mm]}},
		rot/.style={font=\tiny, align=center, inner sep=1pt, text width=2.0cm}
		]
		% Processos (coluna central)
		\node[processo] (p1) at (5.2,0) {1. Curar e inventariar};
		\node[processo] (p2) at (5.2,-2.6) {2. Agrupar e particionar};
		\node[processo] (p3) at (5.2,-5.2) {3. Pré-processar, treinar e selecionar as redes};
		\node[processo] (p4) at (5.2,-7.8) {4. Definir limiares e treinar linhas de base};
		\node[processo] (p5) at (5.2,-10.4) {5. Avaliar no teste reservado};
		% Entidades externas (coluna esquerda)
		\node[entidade] (repo) at (0,0) {Repositório público do 32 CSD v2};
		\node[entidade] (rev) at (0,-1.75) {Revisores (inspeção manual)};
		\node[entidade] (tv) at (0,-5.2) {TorchVision (pesos ImageNet)};
		% Depósitos de dados (coluna direita)
		\node[deposito] (d1) at (10.4,0) {D1\enspace Inventário};
		\node[deposito] (d2) at (10.4,-2.6) {D2\enspace Listas de partição e de dobras};
		\node[deposito] (d3) at (10.4,-5.2) {D3\enspace Modelos e configurações};
		\node[deposito] (d4) at (5.2,-12.7) {D4\enspace Predições e métricas};
		% Fluxos
		\draw[fluxo] (repo.east |- p1.170) -- node[rot, above] {imagens e rótulos originais} (p1.170);
		\draw[fluxo] (rev.east) -- node[rot, below, pos=0.45] {decisões de inclusão} (p1.190);
		\draw[fluxo] (rev.south east) -- node[rot, below left, pos=0.55] {duplicatas confirmadas} (p2.west);
		\draw[fluxo] (p1) -- node[rot, above] {SHA-256, rótulos e procedência} (d1);
		\draw[fluxo] (d1.south west) -- node[rot, sloped, above] {fotografias elegíveis} (p2.north east);
		\draw[fluxo] (p2) -- node[rot, above] {grupos por partição e por dobra} (d2);
		\draw[fluxo] (d2.south west) -- node[rot, sloped, above] {dobras do desenvolvimento} (p3.north east);
		\draw[fluxo] (tv) -- node[rot, above] {pesos pré-treinados} (p3);
		\draw[fluxo] (p3) -- node[rot, above] {configuração escolhida e pesos} (d3);
		\draw[fluxo] (d3.south west) -- node[rot, sloped, above] {modelos retreinados} (p4.north east);
		\draw[fluxo] (d2.east) -- ++(0.7,0) |- node[rot, above, pos=0.75] {partição de limiar e metadados} (p4.east);
		\draw[fluxo] (d2.20) -- ++(1.5,0) |- node[rot, above, pos=0.75] {teste reservado} (p5.east);
		\draw[fluxo] (p4) -- node[rot, right] {limiares e linhas de base fixados} (p5);
		\draw[fluxo] (p5) -- node[rot, right] {predições, métricas e intervalos} (d4);
	\end{tikzpicture}
	\par\vspace{2pt}{\scriptsize Notação: retângulos sombreados são entidades externas; retângulos arredondados, processos; retângulos abertos, depósitos de dados; setas, fluxos de dados.}
	\fonte{Autores.}
\end{figure}

\begin{figure}[H]
	\centering
	\caption{Diagrama de atividades do estudo}
	\label{fig:atividades-estudo}
	\begin{tikzpicture}[
		font=\scriptsize,
		node distance=4.5mm,
		acao/.style={draw, rounded corners=7pt, align=center, execute at begin node={\hyphenpenalty=10000}, minimum height=7mm, text width=5.4cm, inner sep=3pt},
		acaop/.style={draw, rounded corners=7pt, align=center, execute at begin node={\hyphenpenalty=10000}, minimum height=7mm, text width=4.0cm, inner sep=3pt},
		decisao/.style={draw, diamond, aspect=2.6, inner sep=1pt, align=center, font=\tiny},
		barra/.style={fill=black, minimum width=9.6cm, minimum height=1.2mm, inner sep=0pt},
		inicio/.style={circle, fill=black, minimum size=3.2mm, inner sep=0pt},
		fim/.style={circle, draw, double, double distance=1pt, fill=black, minimum size=3mm, inner sep=0pt},
		fluxo/.style={-{Latex[length=1.8mm]}},
		rot/.style={font=\tiny, inner sep=1pt}
		]
		\node[inicio] (ini) {};
		\node[acao, below=of ini] (a1) {Obter o 32 CSD v2 e registrar versão e SHA-256};
		\node[acao, below=of a1] (a2) {Selecionar fotografias clínicas e harmonizar rótulos};
		\node[acao, below=of a2] (a3) {Agrupar duplicatas e fotografias relacionadas};
		\node[decisao, below=of a3] (d1) {Fonte concentra\\uma classe?};
		\node[acaop, right=7mm of d1] (a3b) {Registrar a associação e\\verificar fontes com as duas classes};
		\node[acao, below=6mm of d1] (a4) {Dividir por grupos: desenvolvimento, limiar e teste reservado};
		\node[barra, below=of a4] (fork) {};
		\node[acaop, below=5mm of fork, xshift=-2.4cm] (r1) {ResNet-50: validação cruzada\\em 5 dobras por grupos e\\busca em grade};
		\node[acaop, below=5mm of fork, xshift=2.4cm] (m1) {MobileNetV2: validação cruzada\\em 5 dobras por grupos e\\busca em grade};
		\node[acaop, below=of r1] (r2) {Retreinar com a melhor\\configuração (3 sementes)};
		\node[acaop, below=of m1] (m2) {Retreinar com a melhor\\configuração (3 sementes)};
		\node[barra, below=5mm of r2, xshift=2.4cm] (join) {};
		\node[acao, below=of join] (a5) {Definir limiares na partição de limiar e treinar as linhas de base};
		\node[decisao, below=of a5] (d2) {Modelos e limiares\\fixados e registrados?};
		\node[acao, below=6mm of d2] (a6) {Abrir o teste reservado e calcular métricas e intervalos};
		\node[acao, below=of a6] (a7) {Executar as análises exploratórias (calibração, predição seletiva, ITA)};
		\node[fim, below=of a7] (fimn) {};
		\draw[fluxo] (ini) -- (a1);
		\draw[fluxo] (a1) -- (a2);
		\draw[fluxo] (a2) -- (a3);
		\draw[fluxo] (a3) -- (d1);
		\draw[fluxo] (d1) -- node[rot, above] {[sim]} (a3b);
		\draw[fluxo] (a3b) |- (a4.east);
		\draw[fluxo] (d1) -- node[rot, right] {[não]} (a4);
		\draw[fluxo] (a4) -- (fork);
		\draw[fluxo] (fork.south -| r1.north) -- (r1);
		\draw[fluxo] (fork.south -| m1.north) -- (m1);
		\draw[fluxo] (r1) -- (r2);
		\draw[fluxo] (m1) -- (m2);
		\draw[fluxo] (r2) -- (r2.south |- join.north);
		\draw[fluxo] (m2) -- (m2.south |- join.north);
		\draw[fluxo] (join) -- (a5);
		\draw[fluxo] (a5) -- (d2);
		\draw[fluxo] (d2) -- node[rot, right] {[sim]} (a6);
		\draw[fluxo] (d2.west) -- node[rot, above] {[não]} ++(-1.6,0) |- (a5.west);
		\draw[fluxo] (a6) -- (a7);
		\draw[fluxo] (a7) -- (fimn);
	\end{tikzpicture}
	\par\vspace{2pt}{\scriptsize Notação UML: círculo cheio, início; círculo duplo, fim; losango, decisão; barras, bifurcação e junção de atividades paralelas.}
	\fonte{Autores.}
\end{figure}

\subsection{Pré-processamento das imagens e arquiteturas}
\label{subsec:preprocessamento-arquiteturas}

A ResNet-50 emprega blocos residuais \cite{he2016resnet} e foi usada no
ramo visual do AI4Leprosy por \citeonline{barbieri2022}. A MobileNetV2 usa
blocos com resíduos invertidos e gargalos lineares
\cite{sandler2018mobilenetv2} e foi empregada na SkinLesionNet para
classificar quatro categorias \cite{jaikishore2021}. A escolha reúne uma
rede de referência na literatura revisada e uma rede leve, adequada a
dispositivos com poucos recursos, como os da atenção primária. Usar duas
arquiteturas de famílias diferentes, ao lado de uma regressão logística
como linha de base, atende ao atributo desejável de processar os dados com
mais de um tipo de modelo, já que a conclusão sobre um classificador pode
mudar com a técnica e com os dados \cite{ghotra2015}.

As duas redes serão inicializadas com os pesos \texttt{IMAGENET1K\_V1} do
TorchVision e receberão o mesmo preparo das imagens. Cada fotografia será
redimensionada preservando a proporção, com preenchimento das bordas até
uma resolução quadrada fixa, e normalizada com os valores documentados
para esses pesos. A resolução de partida será de $224 \times 224$ pixels,
a mesma do pré-treinamento; uma resolução maior só será adotada se um
piloto no desenvolvimento mostrar ganho na validação cruzada e for viável
no equipamento, e valerá para as duas redes. A resolução fixa também evita
que o tamanho original do arquivo, associado à fonte, seja visto pela rede
como pista para a classe. O preparo será determinístico na validação e no
teste.

A extração de características é feita pelas camadas convolucionais
pré-treinadas, sem atributos manuais. A agregação por média global
converterá os mapas de características em um vetor de comprimento fixo
por fotografia, e a saída original será substituída por uma camada linear
com uma unidade, que produzirá um \textit{logit}. O treinamento
minimizará a entropia cruzada binária, ponderada pela frequência das
classes na parcela de treino. Na inferência, a função sigmoide produzirá
um escore entre zero e um, convertido em classe pela aplicação de um
limiar.

As transformações de aumento de dados serão geométricas (inversão
horizontal, rotação de até 15 graus e recorte aleatório com reescala) e
fotométricas leves (variação de brilho e contraste). Elas, a ponderação de
classes e quaisquer parâmetros estimados a partir das imagens serão
aplicados ou calculados apenas na parcela de treino de cada dobra, nunca na
validação, na partição de limiar ou no teste. O padrão classifica como
antipadrão o pré-processamento que altera também os dados de teste, como
o rebalanceamento de classes aplicado a eles \cite{acmSigsoftDataScience}.

\subsection{Procedimento experimental e validação cruzada}
\label{subsec:procedimento-validacao}

O treinamento será por transferência: partindo dos pesos pré-treinados, a
camada de saída e as camadas liberadas para atualização serão ajustadas às
fotografias do desenvolvimento (ajuste fino). A validação cruzada servirá
para escolher a configuração de cada rede, e o teste reservado, usado uma
única vez, para estimar o desempenho das configurações escolhidas.

\begin{description}
	\item[Construção das dobras.] A parte do desenvolvimento destinada à
	      validação cruzada será dividida em $k = 5$ dobras por
	      validação cruzada estratificada por grupos: cada grupo de
	      fotografias relacionadas fica inteiro em uma única dobra, e a
	      proporção de grupos positivos é aproximadamente a mesma em todas.
	      Cada grupo é usado para validação exatamente uma vez. As dobras
	      serão geradas uma única vez, com semente fixa, e reutilizadas por
	      todas as configurações das duas redes, o que torna pareadas as
	      comparações entre configurações e entre redes. Se alguma dobra
	      ficar com menos de 20 grupos positivos, $k$ será reduzido para 3
	      antes de qualquer treinamento.
	\item[Treino e validação em cada dobra.] Para cada configuração, a rede
	      é treinada em quatro dobras e avaliada na dobra restante, em
	      rodízio, até que cada dobra tenha servido de validação. A medida de
	      seleção é a área sob a curva ROC (\metrica{M10}) média nas cinco
	      dobras: ela independe do limiar e evita fixar limiares provisórios
	      em cada dobra. O treino em cada dobra é interrompido quando essa
	      área não melhora por cinco épocas consecutivas, com limite de 30
	      épocas, e a época de interrupção é registrada.
	\item[Retreino.] Cada rede é retreinada com a configuração escolhida em
	      todos os grupos da validação cruzada, pelo número de épocas igual à
	      mediana das épocas de interrupção observadas nas dobras. O retreino
	      é repetido com três sementes, e o escore final de cada fotografia é
	      a média dos escores das três execuções. As métricas de cada semente
	      também são relatadas, para mostrar a variação entre execuções.
	\item[Limiares.] O limiar principal é definido na partição reservada para
	      esse fim, com os escores médios do modelo retreinado, maximizando a
	      acurácia balanceada (\metrica{M11}), o que equivale a maximizar o
	      índice de Youden. Em caso de empate, prevalece o limiar com maior
	      sensibilidade (\metrica{M1}). Como ponto de operação secundário,
	      também pré-definido, é registrado o limiar que alcança
	      sensibilidade de pelo menos 0,90 na mesma partição, pela relevância
	      de não perder casos na triagem.
	\item[Abertura do teste.] Os classificadores, os limiares e as linhas de
	      base são fixados e registrados antes da abertura do teste, que é
	      avaliado uma única vez.
\end{description}

A validação cruzada foi preferida a uma única divisão de validação porque
a revisão mostrou que 14 dos 16 estudos que desenvolvem modelos avaliam o
desempenho em uma única divisão (Subseção~\ref{subsec:rq3-metodos}), o que
não informa a variação entre partições. O padrão também lista como
desejável repetir o processamento com diferentes exemplos de treino e
teste, por exemplo por validação cruzada \cite{acmSigsoftDataScience}. A
validação cruzada aninhada, que repetiria a seleção dentro de cada dobra
externa, não será adotada pelo custo computacional; o teste reservado, que
não participa de nenhuma escolha, cumpre o papel de estimativa
independente da seleção.

\begin{algorithm}[H]
	\caption{Procedimento experimental para cada rede}
	\label{alg:procedimento-experimental}
	\footnotesize
	\begin{algorithmic}[1]
		\Require coleção curada $C$ com grupos; grade de hiperparâmetros $\mathcal{H}$; $k = 5$; sementes $S = \{s_1, s_2, s_3\}$
		\Ensure predições, métricas e intervalos de 95\% no teste reservado
		\State $(D, T) \gets$ \Call{DivisãoPorGrupos}{$C$, estratificada, teste de 20\% a 25\%} \Comment{$T$ permanece fechado}
		\State $(D_{cv}, D_{lim}) \gets$ \Call{DivisãoPorGrupos}{$D$, estratificada, 20\% para limiar}
		\State $F_1, \dots, F_k \gets$ \Call{DobrasPorGrupos}{$D_{cv}$, $k$} \Comment{mesmas dobras para as duas redes}
		\ForAll{rede $a \in \{\mbox{ResNet-50}, \mbox{MobileNetV2}\}$}
		\ForAll{configuração $h \in \mathcal{H}$}
		\For{$i = 1$ \textbf{até} $k$}
		\State $M \gets$ \Call{AjusteFino}{$a$, $h$, $D_{cv} \setminus F_i$} \Comment{aumento e pesos de classe só no treino}
		\State $\mathit{AUC}_{h,i} \gets$ \Call{AUC}{$M$, $F_i$};\; $e_{h,i} \gets$ época de interrupção
		\EndFor
		\EndFor
		\State $h^{*}_{a} \gets \arg\max_{h} \frac{1}{k}\sum_{i} \mathit{AUC}_{h,i}$
		\ForAll{semente $s \in S$}
		\State $M_{a,s} \gets$ \Call{AjusteFino}{$a$, $h^{*}_{a}$, $D_{cv}$, épocas $= \mathrm{mediana}_i(e_{h^{*}_{a},i})$, $s$}
		\EndFor
		\State $\bar{p}_a(x) \gets \frac{1}{|S|}\sum_{s} M_{a,s}(x)$ \Comment{escore médio das sementes}
		\State $\tau_a \gets$ limiar que maximiza a acurácia balanceada de $\bar{p}_a$ em $D_{lim}$
		\State $\tau'_a \gets$ limiar com sensibilidade $\geq 0{,}90$ em $D_{lim}$
		\EndFor
		\State Treinar as linhas de base em $D$; fixar e registrar modelos e limiares
		\State Abrir $T$; calcular as métricas e os intervalos por \textit{bootstrap} de grupos ($B = 2.000$)
	\end{algorithmic}
	\fonte{Autores.}
\end{algorithm}

\subsection{Seleção de hiperparâmetros}
\label{subsec:hiperparametros}

O padrão \textit{Data Science} lista entre os atributos desejáveis a
seleção criteriosa dos hiperparâmetros dos modelos, seja a partir de
trabalhos relacionados, seja por um otimizador automático, como a busca em
grade \cite{acmSigsoftDataScience}. A seleção é pertinente neste estudo
por dois motivos. Primeiro, no ajuste fino, a taxa de aprendizado e a
quantidade de camadas liberadas alteram o quanto a rede se afasta dos pesos
pré-treinados, e valores inadequados podem subajustar ou memorizar uma
coleção pequena. Segundo, comparar as duas redes com valores escolhidos
para apenas uma delas favoreceria essa rede; dar a ambas o mesmo espaço de
busca e o mesmo orçamento reduz esse viés.

Para conter o custo, apenas três hiperparâmetros serão buscados, em grade:
a taxa de aprendizado ($10^{-4}$, $3 \times 10^{-4}$ e $10^{-3}$), o
decaimento de pesos ($10^{-4}$ e $10^{-2}$) e as camadas liberadas (apenas
a camada de saída, o último bloco convolucional ou toda a rede), em 18
configurações por rede. Os demais serão fixados com base nas
recomendações dos pesos do TorchVision e nos estudos revisados: otimizador
AdamW, lotes de 32 fotografias, resolução de $224 \times 224$ pixels e o
conjunto de aumentos da Subseção~\ref{subsec:preprocessamento-arquiteturas}.
A busca usará somente as dobras da validação cruzada; a partição de limiar
e o teste não participam. A grade, a configuração escolhida e a área sob a
curva ROC de cada configuração em cada dobra serão relatadas.

\subsection{Linhas de base e comparação com o estado da arte}
\label{subsec:linhas-base}

O padrão exige uma de três alternativas: discutir as linhas de base do
estado da arte, com seus pontos fortes, fracos e limitações; explicar por
que elas não existem; ou argumentar de forma convincente que a comparação
direta é impraticável \cite{acmSigsoftDataScience}. O planejamento trata
as três, porque cada uma se aplica a uma parte do problema.

\begin{description}
	\item[Linhas de base do estado da arte.] Entre os estudos revisados, três
	      abordagens são as referências mais próximas. O AI4Leprosy
	      \cite{barbieri2022} usa a ResNet-50 no ramo visual e combina
	      imagens com dados clínicos; seu ponto forte é a amostra brasileira
	      com registro por paciente, e suas limitações, para este estudo, são
	      o acesso restrito e a dependência de dados clínicos que o 32 CSD não
	      contém. O classificador do aplicativo da OMS para doenças tropicais
	      negligenciadas de pele \cite{pons2026} tem a seu favor a validação
	      externa independente de \citeonline{deps2026}, mas não está
	      disponível para execução local. A SkinLesionNet
	      \cite{jaikishore2021} usa a MobileNetV2, o que mostra a viabilidade
	      de uma rede leve, mas classifica quatro categorias com partições por
	      imagem. Essas abordagens orientam a escolha das duas redes, sem
	      servir de comparação numérica.
	\item[Ausência de linha de base para a tarefa definida.] Nenhum dos
	      estudos incluídos usa o 32 CSD (Apêndice~\ref{apend:inventario-datasets})
	      nem relata a distinção binária entre hanseníase e outras condições
	      de pele com partições por grupos de fotografias relacionadas. Não
	      há, portanto, um resultado do estado da arte para a mesma tarefa,
	      nos mesmos dados e com o mesmo controle de vazamento.
	\item[Impraticabilidade da comparação direta.] Mesmo os valores
	      publicados não são comparáveis: as coleções, as classes e as
	      partições diferem, apenas 2 dos 16 estudos que desenvolvem modelos
	      declaram separar os dados por paciente
	      (Subseção~\ref{subsec:rq3-metodos}), e há sobreposição de origem
	      entre o 32 CSD e as imagens usadas por \citeonline{eita2025} e por
	      \citeonline{lima2024skin}, o que viola a independência entre
	      dados de comparação. Os valores da literatura serão usados apenas
	      para discussão qualitativa (questão Q3).
\end{description}

Na falta de uma linha de base externa comparável, duas linhas de base
internas serão treinadas no desenvolvimento e avaliadas no mesmo teste. A
primeira é um classificador trivial, que atribui sempre a classe
majoritária ou sorteia classes na proporção do treino, cuja acurácia
balanceada (\metrica{M11}) esperada é de 0,50 e serve de piso. A segunda é
uma regressão logística que usa apenas atributos do arquivo, como
dimensões, proporção e tamanho, sem olhar o conteúdo da imagem. Se essa
regressão superar o acaso, os atributos de origem carregam informação
sobre a classe, e o ganho das redes deve ser interpretado em relação a ela
(questão Q6).

\section{Instrumentação}
\label{sec:instrumentacao}

Seguindo \citeonline{wohlin2024}, a instrumentação reúne os objetos do
estudo, as diretrizes que orientam sua execução e os instrumentos de
medição. Como o estudo é computacional, inclui também a infraestrutura de
hardware e software, exigida pelo padrão \textit{Data Science}
\cite{acmSigsoftDataScience}.

\begin{description}
	\item[Objetos.] A versão 2 do 32 CSD, identificada pelo nome e pelo
	      SHA-256 de cada arquivo, e os pesos \texttt{IMAGENET1K\_V1} da
	      ResNet-50 e da MobileNetV2 distribuídos pelo TorchVision.
	\item[Diretrizes.] O protocolo de curadoria
	      (Subseção~\ref{subsec:curadoria-procedencia}), com os critérios de
	      seleção, as regras de harmonização dos rótulos e as regras de
	      agrupamento; um guia de inspeção manual para os dois revisores, com
	      as categorias de decisão (incluir, excluir por modalidade, excluir
	      por rótulo inconclusivo, agrupar) e o registro das divergências; e
	      este plano, com as decisões fixadas antes da abertura do teste.
	\item[Instrumentos de medição.] O inventário em formato tabular, com um
	      registro por arquivo; os \textit{scripts} de curadoria, que calculam
	      os hashes criptográficos e perceptuais e propõem os pares candidatos
	      a duplicata; os \textit{scripts} de particionamento, que geram as
	      listas de grupos por partição e por dobra; os \textit{scripts} de
	      treinamento e de inferência, que registram configuração, semente,
	      épocas e escores por fotografia; e os \textit{scripts} de avaliação,
	      que calculam as métricas \metrica{M1} a \metrica{M6} e \metrica{M10}
	      a \metrica{M12}, a matriz de confusão e os intervalos por
	      \textit{bootstrap} de grupos, além da medição do ITA nas regiões de
	      pele íntegra. Os \textit{scripts} de avaliação serão verificados com
	      casos de teste de resultado conhecido antes do uso.
	\item[Registro dos experimentos.] Cada execução gravará um arquivo de
	      configuração versionado, com a lista de partição usada, os
	      hiperparâmetros, a semente, as versões das bibliotecas e o
	      identificador do código no controle de versão, de modo que cada
	      resultado possa ser associado à execução que o produziu.
	\item[Infraestrutura.]
	      \phantomsection\label{subsec:infraestrutura}
	      A implementação usará Python, PyTorch e TorchVision, com
	      scikit-learn para as partições, a regressão logística e parte das
	      métricas. Serão relatados o processador e a placa gráfica, a
	      memória disponível, o sistema operacional, as versões das
	      bibliotecas e o tempo de treino e de inferência por fotografia em
	      cada rede, informação relevante para a MobileNetV2 como alternativa
	      leve. As sementes serão fixadas, e as operações não determinísticas
	      da placa gráfica serão documentadas.
	\item[Pacote de replicação.] Serão disponibilizados, em repositório
	      público, o inventário, as listas de fotografias por partição e por
	      dobra, as regras de seleção, os arquivos de configuração e os
	      \textit{scripts} de curadoria, de treinamento e de avaliação, o que
	      permite reexecutar o estudo a partir da versão pública da coleção
	      \cite{acmSigsoftDataScience,wohlin2024}.
\end{description}

\section{Medidas de desempenho e análises complementares}
\label{sec:medidas-avaliacao-empirica}

A medida principal será a acurácia balanceada (\metrica{M11}), média entre
a sensibilidade para hanseníase (\metrica{M1}) e a especificidade para as
outras condições (\metrica{M2}). Ela atribui o mesmo peso às duas classes
e não é inflada pelo desequilíbrio da coleção, em que a hanseníase
corresponde a cerca de 13\% das imagens. Serão relatadas também a
sensibilidade, com destaque pela finalidade de triagem, a especificidade,
a área sob a curva ROC (\metrica{M10}), a matriz de confusão e as
contagens, para os dois classificadores e as duas linhas de base. A
acurácia (\metrica{M5}), a precisão (\metrica{M3}) e o $F_1$ (\metrica{M4})
também serão relatados, para situar os resultados em relação à literatura,
com a ressalva de que dependem da prevalência das classes no teste.

Os intervalos de confiança de 95\% serão calculados por
\textit{bootstrap} de grupos do teste, com 2.000 reamostragens
estratificadas por classe e intervalos por percentis. Reamostrar grupos,
e não fotografias, preserva a dependência entre imagens relacionadas.
Como os dois modelos serão avaliados nas mesmas fotografias, a diferença
entre suas acurácias balanceadas será relatada com intervalo pareado, em
que cada reamostragem usa os mesmos grupos para os dois modelos; a
diferença e seu intervalo expressam o tamanho do efeito, sem teste de
hipótese adicional, o que evita o antipadrão de relatar significância sem
tamanho de efeito \cite{acmSigsoftDataScience}. Também serão relatados a
média e o desvio padrão da área sob a curva ROC entre as dobras da
validação cruzada e a variação entre as três sementes do retreino final.

A calibração dos escores será examinada pelo escore de Brier e pelo erro
esperado de calibração (\metrica{M12}) e pelo diagrama de confiabilidade,
com intervalos de agrupamento definidos no desenvolvimento
\cite{guo2017calibration}. A análise exploratória de predição seletiva
examinará a relação entre cobertura e acurácia balanceada com regras de
abstenção definidas no desenvolvimento, informando a fração de fotografias
classificadas e os denominadores por classe. Falsos positivos e falsos
negativos serão inspecionados por categoria original e, quando possível,
por fonte.

Em resposta à ausência de resultados por tom de pele observada na revisão,
será considerada uma análise exploratória da diferença de desempenho entre
grupos (\metrica{M6}) por faixas do ângulo de tipologia individual (ITA)
\cite{chardon1991ita}, se houver regiões verificáveis de pele íntegra. A
região de medição será inspecionada para excluir fundo e objetos, e o ITA
será interpretado como medida colorimétrica nas condições de captura,
sujeita a iluminação e câmera. As métricas por faixa só serão estimadas
quando houver ao menos 10 grupos positivos e 10 negativos na faixa; nos
demais casos, serão informados apenas os denominadores e as exclusões.

\section{Pressupostos técnicos e ameaças à validade}
\label{sec:validade-avaliacao-empirica}

A interpretação seguirá as quatro dimensões da validade discutidas por
\citeonline{wohlin2024}, precedidas pelos pressupostos técnicos próprios
de estudos de ciência de dados, que o padrão pede que sejam explicitados
\cite{acmSigsoftDataScience}. Para cada ameaça, indica-se a medida prevista
para reduzi-la ou o limite que ela impõe à interpretação.

\subsection{Pressupostos técnicos}
\label{subsec:pressupostos-tecnicos}

\begin{enumerate}
	\item \textbf{Independência entre grupos.} As estimativas e os
	      intervalos supõem que grupos diferentes são independentes. Como a
	      coleção não identifica pacientes, fotografias da mesma pessoa que
	      não sejam detectadas como relacionadas podem cair em grupos
	      distintos. O agrupamento por similaridade reduz, mas não elimina,
	      essa dependência residual.
	\item \textbf{Rótulos como referência.} O estudo supõe que os rótulos da
	      coleção representam a condição fotografada. Erros de rótulo afetam
	      treino e teste da mesma forma e não são detectáveis apenas pelos
	      dados; a inspeção manual de amostras estima sua ocorrência.
	\item \textbf{Prevalência da coleção.} A proporção de hanseníase no
	      teste, cerca de 13\%, decorre da composição da coleção e não da
	      prevalência em um serviço de saúde. Medidas dependentes da
	      prevalência, como precisão, $F_1$ e acurácia, e os limiares
	      definidos na partição de limiar não se transferem diretamente para
	      outro contexto.
	\item \textbf{Adequação do pré-treinamento.} O ajuste fino supõe que as
	      representações aprendidas no ImageNet são úteis para fotografias de
	      pele. A comparação com as linhas de base mostra se as redes de fato
	      extraem informação do conteúdo das imagens.
\end{enumerate}

\subsection{Validade de construto}

\begin{itemize}
	\item \textbf{Referência sem confirmação clínica.} O desempenho expressa a
	      concordância com os rótulos da coleção, que não foram confirmados
	      por exame clínico, baciloscopia ou biópsia neste estudo. Os
	      resultados serão descritos como concordância com a anotação da
	      coleção, e não como acurácia diagnóstica.
	\item \textbf{Definição das classes.} O agrupamento das três formas de
	      hanseníase em uma classe positiva e das demais condições em uma
	      classe negativa delimita o que a medida representa. A
	      sensibilidade será também relatada por categoria original de
	      hanseníase, e os erros, por categoria da classe negativa.
	\item \textbf{Viés de operação única.} O construto ``capacidade de
	      distinguir hanseníase'' é medido em uma só coleção e com uma só
	      modalidade de imagem. Por isso, o estudo não generaliza o construto
	      para outras coleções nem para outras modalidades.
	\item \textbf{Escolha das métricas.} Uma única métrica não descreve o
	      desempenho na triagem. A acurácia balanceada é a medida principal,
	      mas a sensibilidade, a especificidade, a área sob a curva ROC e a
	      calibração serão relatadas em conjunto, com a matriz de confusão.
\end{itemize}

\subsection{Validade interna}

\begin{itemize}
	\item \textbf{Vazamento entre partições.} Duplicatas, recortes e
	      fotografias da mesma lesão em partições diferentes inflariam o
	      desempenho \cite{kapoor2023}. O agrupamento por hash, similaridade
	      perceptual e inspeção visual e a divisão por grupos em todas as
	      etapas, inclusive nas dobras, controlam esse risco.
	\item \textbf{Atalhos ligados à fonte.} Bordas, marcas d'água,
	      resolução ou estilo de captura associados a uma fonte podem
	      permitir prever a classe sem olhar a lesão. A auditoria de
	      procedência, a resolução fixa, a linha de base com metadados e a
	      inspeção dos erros por fonte tratam esse risco.
	\item \textbf{Sobreposição com o pré-treinamento.} Os pesos do ImageNet e
	      as imagens públicas da coleção podem ter origem comum. A procedência
	      dos pesos será examinada quando suas informações permitirem.
	\item \textbf{Comparação desigual entre redes.} Uma rede poderia ser
	      favorecida por um espaço de busca ou um orçamento maiores. As duas
	      redes terão a mesma grade, as mesmas dobras, o mesmo preparo das
	      imagens e o mesmo número de sementes.
	\item \textbf{Uso indevido do teste.} Consultar o teste durante o
	      desenvolvimento tornaria a estimativa otimista. O teste ficará
	      fechado até que modelos, limiares e linhas de base sejam fixados e
	      registrados, e será avaliado uma única vez.
	\item \textbf{Erros de instrumentação.} Falhas nos \textit{scripts}
	      poderiam produzir métricas incorretas. Os \textit{scripts} de
	      avaliação serão verificados com casos de resultado conhecido, e o
	      código será revisado pelo outro autor antes da abertura do teste.
	\item \textbf{Aleatoriedade do treinamento.} A inicialização da camada de
	      saída, a ordem dos lotes e operações não determinísticas da placa
	      gráfica alteram o resultado. As sementes serão fixadas, e o
	      retreino será repetido com três sementes, com a variação relatada.
\end{itemize}

\subsection{Validade de conclusão}

\begin{itemize}
	\item \textbf{Baixo número de grupos positivos.} Com poucos grupos de
	      hanseníase no teste, os intervalos ficam largos e diferenças
	      pequenas entre redes não podem ser distinguidas. A meta de cerca de
	      100 grupos positivos no teste foi definida para limitar a
	      semiamplitude a cerca de 8 pontos percentuais, e os denominadores
	      serão sempre relatados.
	\item \textbf{Comparações múltiplas.} Muitas métricas, análises por faixa
	      e pontos de operação aumentam a chance de achados ao acaso. A
	      acurácia balanceada é definida de antemão como medida principal e
	      a diferença pareada entre as redes como comparação principal; as
	      análises de calibração, de predição seletiva e por faixa de ITA são
	      exploratórias e não sustentam conclusões confirmatórias.
	\item \textbf{Pressupostos do \textit{bootstrap}.} O \textit{bootstrap}
	      por percentis pode subestimar a incerteza com poucos grupos ou
	      métricas próximas de zero ou de um. Nesses casos, os intervalos
	      serão interpretados com cautela e as contagens absolutas serão
	      apresentadas.
	\item \textbf{Confiabilidade das medidas.} O ITA depende da iluminação e
	      da câmera, e os rótulos podem conter erros. Por isso, as faixas de
	      ITA são tratadas como medida colorimétrica aproximada, e as análises
	      por faixa só serão feitas com ao menos 10 grupos de cada classe.
\end{itemize}

\subsection{Validade externa}

\begin{itemize}
	\item \textbf{Coleção única.} O teste será interno à coleção pública. A
	      combinação de fontes, a composição das categorias e as condições de
	      captura delimitam o conjunto de fotografias ao qual os resultados
	      se aplicam.
	\item \textbf{Condições de captura.} Imagens de atlas e de coleções
	      curadas tendem a ter enquadramento e iluminação melhores que
	      fotografias de celular feitas em serviços de atenção primária. O
	      desempenho pode cair nessas condições, que não são avaliadas.
	\item \textbf{População representada.} A distribuição de tons de pele,
	      idades e formas clínicas na coleção não é documentada. Se a
	      análise por ITA for viável, ela apenas descreverá a variação nas
	      faixas com amostra suficiente.
	\item \textbf{Contexto de uso.} A transferência para novas coleções e
	      para o atendimento no SUS é uma questão de avaliação posterior,
	      que exige validação externa prospectiva. Os resultados não serão
	      apresentados como evidência de efetividade clínica.
\end{itemize}

\clearpage
\endgroup
